% =====================================================================
%  5. Implementação do esgotamento sanitário
%  (Atividades 3.2.2.4 e 3.2.2.5 do MOP)
% =====================================================================
\section{Implementação dos sistemas de esgotamento sanitário}
\label{sec:esgotamento}

O esgotamento sanitário será resolvido por sistemas descentralizados, isto é,
soluções individuais instaladas em cada domicílio, adequadas à baixa
densidade das áreas rurais. A implementação tem duas portas de entrada
(\autoref{fig:esgotamento-geral}). A frente das comunidades atende 70\% dos
domicílios das comunidades que recebem o sistema coletivo de abastecimento
de água, o que corresponde a 4.067 domicílios (\atividade{3.2.2.4}). A
frente dos mananciais atende 2.300 domicílios em propriedades rurais
situadas nas áreas prioritárias do PSH-PR em que atua o IDR-Paraná
(\atividade{3.2.2.5}). As duas frentes diferem na forma como os domicílios
são identificados e contratados, mas, a partir do projeto de cada
propriedade, seguem o mesmo procedimento de instalação, verificação,
capacitação, medição e pagamento \cite{mop2026}.

\begin{figure}[htbp]
  \centering
  \caption{Implementação do esgotamento sanitário descentralizado, do começo ao fim}
  \label{fig:esgotamento-geral}
  \fluxograma{fluxograma-esgotamento-geral}
  \fonte{Elaborado pelo IAT (2026) com base em \citeonline{mop2026} e no POA 2026--2027.}
\end{figure}

% ---------------------------------------------------------------------
\subsection{Identificação dos domicílios e definição das soluções}
\label{subsec:esgotamento-demanda}

Na frente das comunidades, o ponto de partida é o trabalho técnico-social
descrito na \autoref{subsec:tts}. Depois da seleção e da mobilização das
prefeituras e das comunidades, o diagnóstico domiciliar identificará como
cada família dispõe hoje o seu esgoto. É a partir dessa informação que se
definem os domicílios a serem atendidos \cite{mop2026}. Antes de contratar a
instalação, o IAT, com o IDR-Paraná, fará um estudo para definir as melhores
alternativas de sistemas de esgotamento sanitário para comunidades rurais.
Esse estudo será conduzido pela equipe interna entre outubro e dezembro de
2026, segundo o POA, e servirá de base técnica para os termos de referência.

Na frente dos mananciais, a demanda nasce dos Planos Integrados de Propriedade
(PIP) elaborados pelo IDR-Paraná no âmbito do Subcomponente 2.3. Concluídos os
PIP, será iniciado o processo de instalação dos sistemas nas propriedades
situadas nos mananciais prioritários do PSH-PR \cite{mop2026}.

\pendencia{O MOP não indica a tecnologia de tratamento prevista nem as normas
técnicas de projeto. Depois do estudo de alternativas, citar as normas
aplicáveis, como a ABNT NBR 7229 (tanques sépticos) e a ABNT NBR 13969
(unidades de tratamento complementar e disposição final).}

% ---------------------------------------------------------------------
\subsection{Contratação}
\label{subsec:esgotamento-contratacao}

Em ambas as frentes, o IAT elaborará os termos de referência para a
aquisição das soluções descentralizadas e para a contratação dos serviços de
obras civis de instalação. Aprovados os TR, serão abertos os editais,
realizadas as licitações e contratadas as empresas \cite{mop2026}. O POA
distingue duas contratações. Na frente das comunidades, a aquisição de
material para implantação do esgotamento terá TR entre janeiro e março de
2027 e licitação de abril a dezembro de 2027. Na frente do IDR-Paraná, a
aquisição de materiais e a implantação dos sistemas serão contratadas em
conjunto, com TR entre abril e junho de 2027 e licitação de julho a dezembro de
2027. A supervisão dessas obras ficará a cargo da mesma empresa contratada
para supervisionar as obras de abastecimento (\autoref{subsec:obras}).

\pendencia{O POA fala em ``rede de esgotamento sanitário'' na contratação da
frente das comunidades, enquanto o MOP prevê sistemas descentralizados
individuais. Uniformizar a terminologia.}

% ---------------------------------------------------------------------
\subsection{Projeto, instalação e verificação técnica}
\label{subsec:esgotamento-instalacao}

Contratada a empresa, ela elaborará um projeto técnico específico para cada
residência, considerando as características do terreno, o número de moradores
e a solução de tratamento mais adequada. Em seguida, executará as obras civis
necessárias à instalação das estruturas e do dispositivo de esgotamento
domiciliar. Concomitantemente à instalação, os proprietários serão orientados
sobre o funcionamento do sistema, as práticas adequadas de uso e a manutenção
preventiva, com registro dessa capacitação \cite{mop2026}.

Concluída a instalação, será feita uma verificação técnica para atestar que o
sistema funciona conforme os requisitos definidos e para registrar a sua
garantia. Essa verificação funciona como um ponto de decisão. Se o sistema
atender às condições técnicas, é validado e segue para a medição. Se estiver
em desacordo, a empresa fará as correções necessárias e o sistema será
verificado novamente, até que seja aprovado (\autoref{fig:esgotamento-instalacao}).
Os testes de estanqueidade após a instalação, previstos entre as medidas
socioambientais do Subcomponente~3.2, reduzem o risco de contaminação do solo
e do lençol freático por instalação incorreta.

\begin{figure}[htbp]
  \centering
  \caption{Projeto, instalação e verificação técnica dos sistemas descentralizados}
  \label{fig:esgotamento-instalacao}
  \fluxograma{fluxograma-esgotamento}
  \fonte{Elaborado pelo IAT (2026) com base em \citeonline{mop2026}.}
\end{figure}

% ---------------------------------------------------------------------
\subsection{Medição, comprovação e pagamento}
\label{subsec:esgotamento-medicao}

Como os sistemas são instalados propriedade por propriedade, o pagamento
depende de uma cadeia de comprovação que envolve a empresa, o proprietário e
o IAT (\autoref{fig:esgotamento-medicao}). A empresa registrará, em relatórios
de medição, o andamento das atividades e as intervenções executadas em cada
propriedade, e reunirá a documentação que comprova a elaboração e a execução
dos projetos, com as evidências das intervenções. O proprietário, por sua
vez, assinará o registro que confirma a execução das melhorias em sua
propriedade \cite{mop2026}.

Com esse material, o IAT fará uma auditoria documental e visitas técnicas de
verificação \textit{in loco} para confirmar a execução das intervenções. A
empresa consolidará então as informações no relatório de medição final, e o
IAT avaliará por amostragem, também \textit{in loco}, a qualidade das unidades
atendidas. Só depois dessas verificações serão realizadas a ordenação das
despesas e o pagamento dos serviços, seguidos do encerramento formal do
processo. O relatório técnico aprovado e o comprovante de pagamento
constituem a evidência de conclusão de cada lote \cite{mop2026}.

\begin{figure}[htbp]
  \centering
  \caption{Medição, comprovação, verificação independente e pagamento, por responsável}
  \label{fig:esgotamento-medicao}
  \fluxograma{fluxograma-esgotamento-medicao}
  \fonte{Elaborado pelo IAT (2026) com base em \citeonline{mop2026}.}
\end{figure}

\pendencia{O MOP não define o tamanho da amostra da verificação \textit{in
loco} nem a diferença entre ``verificação \textit{in loco}'' e ``amostragem
\textit{in loco}''. Também não indica quem executa cada verificação; nesta
estrutura, elas foram atribuídas ao IAT, responsável pela validação.
Confirmar.}

% ---------------------------------------------------------------------
\subsection{Particularidades da frente do IDR-Paraná}
\label{subsec:esgotamento-idr}

A frente do IDR-Paraná segue o mesmo procedimento de projeto, instalação,
verificação, capacitação, medição e pagamento descrito acima. As diferenças
estão na origem da demanda, no território e nos responsáveis
(\autoref{qua:diferencas-esgotamento}). Os domicílios são identificados pelos
PIP e não pelo trabalho técnico-social. O território é o dos mananciais de
abastecimento público, onde o esgoto doméstico não tratado representa risco
direto para a água que abastece as cidades. O projeto é chamado de projeto
sanitário doméstico da propriedade. Por fim, a validação é compartilhada entre
o IAT e o IDR-Paraná e não há dependência das prefeituras \cite{mop2026}.

\begin{quadro}[htbp]
\caption{Diferenças entre as frentes de esgotamento sanitário}
\label{qua:diferencas-esgotamento}
\small
\renewcommand{\arraystretch}{1.25}
\begin{tabularx}{\textwidth}{|P{3.2cm}|L|L|}
\hline
\cab{Aspecto} & \cab{Frente comunidades (3.2.2.4)} & \cab{Frente mananciais -- IDR-Paraná (3.2.2.5)} \\ \hline
Origem da demanda & Diagnóstico do trabalho técnico-social &
  Planos Integrados de Propriedade (PIP) do IDR-Paraná \\ \hline
Território & Comunidades atendidas com abastecimento de água &
  Mananciais prioritários do PSH-PR \\ \hline
Meta & 4.067 domicílios (70\% dos atendidos com água) & 2.300 domicílios \\ \hline
Contratação (POA) & Aquisição de material; TR jan--mar/2027 &
  Aquisição de materiais e implantação; TR abr--jun/2027 \\ \hline
Prazo (planilha de custos) & 2028 a 2033 & 2029 a 2033 \\ \hline
Dependências & Prefeituras & -- \\ \hline
Evidência & \multicolumn{2}{l|}{Relatório técnico aprovado e comprovante de pagamento, por lote} \\ \hline
Validação & IAT & IAT e IDR-Paraná \\ \hline
\end{tabularx}
\fonte{Adaptado de \citeonline[Quadros 30 e 31]{mop2026} e do POA 2026--2027.}
\end{quadro}
